\section{Conclusion}\label{sec:conclusion}

Agents ask humans before they act, and the request is treated as a second factor. We tested when it is one. Binding
the approval to the exact effect removed every post-approval attack against \sys{} in our experiments, and AP2 and
Microsoft's ADR~0030 implementation already do it well. Rendering is where approvals still break.
Padding, spoofed labels and displays that do not distinguish two effects succeeded against real designs, and a
hidden blind copy passed every model approver that read the agent's description. When the agent narrated a forged
action, a narrated card let \adFExNarr{} of AgentDojo attacks through, while a trusted card let \adFExAsu{}
through. Language models used as approvers caught \llmCatchNarr{} of attacks from the agent's
description and \llmCatchCard{} from the trusted card. The second device did not change any measured outcome. Its
value lies in the trust domain and the verifiable receipt, which the proofs establish and the experiments do not
test. Flooding, fatigue and effects that policy allows without asking remain open. Our contribution is specific: a set of attacks, an audit and a set of measurements that show which half of agent
MFA existing designs already address and which half remains open. The code, data, models and the scripts that produce every number in this paper are available
at \codeurl{}.
